\chapter{Partial derivatives and the multivariable chain}

\begin{orient}
\textbf{What this chapter is for.} Nothing you learned about differentiation in one variable is repealed when a second variable appears. The rules are the same rules, the table of derivatives is the same table, and the algebra is the same algebra. What changes is the bookkeeping: which symbols are frozen while you differentiate, how many terms a chain rule must have, and which of several possible derivatives the problem is actually asking for. Almost every error in this chapter is a bookkeeping error, not a calculus error, and the techniques below are all devices for making the bookkeeping visible so that it can be checked. You will learn to hold variables fixed deliberately rather than by habit, to draw the dependency tree of a composite and read the chain rule off it, to use the gradient as the single object that answers every question about direction, to differentiate an implicit relation in one line by $-F_{x}/F_{y}$, to propagate small errors through a formula, and to test a pair of functions for exactness, which is the last technique in this part and the first technique in the part on differential equations.
\end{orient}

\section{Holding the others fixed}

A partial derivative is an ordinary derivative taken along a coordinate line. To compute $\pdv{f}{x}$ at a point you restrict $f$ to the horizontal line through that point, which turns it into a function of the single variable $x$, and differentiate that function in the ordinary way. Every rule from the single-variable theory therefore applies verbatim: product, quotient, chain, the derivatives of $\sin$ and $\ln$ and $\arctan$. The only new instruction is that $y$ is now a constant, and the only new failure is forgetting that it is one, or forgetting that it stops being one.

That instruction is easy to state and surprisingly hard to obey, because the mind does not treat all constants alike. Nobody differentiating $7x^{3}$ writes anything but $21x^{2}$, but a large number of readers differentiating $y\,\eu^{xy}$ with respect to $x$ write $y\,\eu^{xy}$ and stop, because the $y$ in the exponent is doing two jobs at once and only one of them registers. The discipline that fixes this is mechanical: before differentiating, rewrite the expression with the frozen variable replaced by a letter that looks like a constant. Differentiating $y\,\eu^{xy}$ becomes differentiating $c\,\eu^{cx}$, which is $c^{2}\eu^{cx}$, so $\pdv{}{x}\bigl(y\,\eu^{xy}\bigr)=y^{2}\eu^{xy}$. The substitution costs five seconds and removes the entire error class.

\tech{Partial differentiation with the others frozen}{core}
\cue Any $f(x,y)$ or $f(x,y,z)$ with $f_{x}$, $f_{xy}$, or $\dfrac{\partial^{n+m}f}{\partial x^{n}\partial y^{m}}$ requested.
\method Differentiate with respect to the named variable, treating every other variable as a numerical constant. For a mixed partial, do the derivatives one at a time in whichever order is cheaper: Clairaut's theorem says $f_{xy}=f_{yx}$ wherever both are continuous, so the order is yours to choose. Two structural facts do most of the work. Any term depending on only one of the variables is annihilated by any mixed partial. And any term whose degree in $x$ is less than $n$ is annihilated by $\partial_{x}^{n}$, no matter what it does in the other variables.

\begin{worked}
$f(x,y)=10x^{4}y^{3}-\eu^{3xy}+1$, evaluated at $(2,0)$.

Take $f_{x}$ with $y$ frozen: $f_{x}=40x^{3}y^{3}-3y\,\eu^{3xy}$. Every term carries a factor of $y$, so $f_{x}(2,0)=0$. Take $f_{y}$ with $x$ frozen: $f_{y}=30x^{4}y^{2}-3x\,\eu^{3xy}$, and at $(2,0)$ the first term dies and the second is $-6$. For the mixed partial, differentiate $f_{x}$ with respect to $y$, remembering that $-3y\,\eu^{3xy}$ needs the product rule:
\[
f_{xy}=120x^{3}y^{2}-3\eu^{3xy}-9xy\,\eu^{3xy},
\]
so $f_{xy}(2,0)=0-3-0=-3$. Differentiating $f_{y}$ with respect to $x$ instead gives $120x^{3}y^{2}-3\eu^{3xy}-9xy\,\eu^{3xy}$ as well, which is the theorem doing its job. Central differences confirm $f_{y}(2,0)=-6.0000$ and $f_{xy}(2,0)=-3.0000$.

Now the same function with an order nobody intends you to reach by hand: $\dfrac{\partial^{9}f}{\partial x^{5}\partial y^{4}}$ at $(2,0)$. Two of the three terms die before any work. The polynomial $10x^{4}y^{3}$ has degree $4$ in $x$ and is annihilated by $\partial_{x}^{5}$, and the constant $1$ is annihilated by everything. Only $-\eu^{3xy}$ is left, and with $x$ the active variable and $y$ frozen it is a pure exponential, so
\[
\partial_{x}^{5}\bigl(-\eu^{3xy}\bigr)=-(3y)^{5}\eu^{3xy}.
\]
Every remaining derivative is in $y$, and Leibniz on the product $(3y)^{5}\cdot\eu^{3xy}$ with only four $y$-derivatives available can lower the exponent of $y^{5}$ to at most $y^{1}$:
\[
\partial_{y}^{4}\Bigl[(3y)^{5}\eu^{3xy}\Bigr]
=\sum_{j=0}^{4}\binom{4}{j}\frac{5!}{(5-j)!}\,3^{5}y^{5-j}\,(3x)^{4-j}\eu^{3xy},
\]
and every term carries a positive power of $y$. At $y=0$ they all vanish, so the ninth mixed partial is exactly $0$ at $(2,0)$, and at every other point of the $x$-axis as well. Finite differences on the lower analogue $\partial_{x}^{3}\partial_{y}^{2}\eu^{3xy}$ at $(2,0)$ return $3.44$ at step $0.05$ and $0.52$ at step $0.02$, shrinking like $h^{2}$ towards the exact zero. The whole computation was two annihilation rules and one count of how many $y$-derivatives were on offer.
\end{worked}

\begin{trap}
Letting a frozen variable participate in the chain rule anyway. Differentiating $x^{zy}$ with respect to $y$ is not a power rule: with $x$ and $z$ frozen the function is $c^{zy}$ in the exponent, so the derivative is $x^{zy}\ln x\cdot z$. The factor $\ln x$ appears in the $y$-partial and in the $z$-partial and does not appear in the $x$-partial, where the power rule genuinely applies and gives $zy\,x^{zy-1}$. A single expression can require three different rules for its three partials.
\end{trap}

\begin{seealsob}The dependency tree; mixed-partial coefficient extraction; the exactness test.\end{seealsob}

Clairaut's theorem is worth more than a footnote, because it is the only licence you have to reorder work, and because the hypothesis it carries is the kind that is usually true and occasionally decisive. The statement is that if $f_{xy}$ and $f_{yx}$ both exist and are continuous on a neighbourhood of a point, they are equal there. In practice this means: choose the order that kills the most terms first, or the order that avoids a product rule.

\begin{keynote}[Where the mixed partials disagree]
Continuity is a real hypothesis. Let
\[
f(x,y)=\frac{xy\,(x^{2}-y^{2})}{x^{2}+y^{2}},\qquad f(0,0)=0 .
\]
Along the $y$-axis, $f(h,y)-f(0,y)=f(h,y)$, so $f_{x}(0,y)=\lim_{h\to0}\tfrac{y(h^{2}-y^{2})}{h^{2}+y^{2}}=-y$; by the antisymmetry of $f$ under swapping $x$ and $y$, $f_{y}(x,0)=x$. Differentiating those one-variable functions gives
\[
f_{xy}(0,0)=\dvop{y}(-y)=-1,\qquad f_{yx}(0,0)=\dvop{x}(x)=+1 .
\]
Numerically, nested difference quotients with an inner step of $10^{-6}$ and an outer step of $10^{-4}$ return $-0.999800$ and $+0.999800$. Both second partials exist at the origin and they differ by $2$. What fails is continuity: $f_{xy}$ has no limit at the origin, since its value depends on the direction of approach. Every function you will meet in a problem is smooth enough for Clairaut, but the reason to know this example is that it tells you what the theorem is actually asserting, which is a statement about the second derivatives near the point and not merely at it.
\end{keynote}

\section{The dependency tree}

The chain rule in one variable has one term because a composite $f(g(x))$ has one path from $x$ to $f$. That is the whole content of the rule, and it generalises without modification: the derivative of a composite is a sum over paths, one term per route from the differentiating variable up to the output, and each term is the product of the derivatives encountered along that route. Draw the routes and the formula writes itself. Refuse to draw them and you will omit one, which is the single commonest error in this chapter and the reason the tree is taught here as a technique rather than as an illustration.

The picture for $f\bigl(x(s,t),y(s,t)\bigr)$ is this.

\begin{center}
\begin{tikzpicture}[node distance=6mm and 8mm]
\node[dnode] (f) {$f$};
\node[dnode, below left=9mm and 9mm of f] (x) {$x$};
\node[dnode, below right=9mm and 9mm of f] (y) {$y$};
\node[dnode, below left=9mm and 1mm of x] (s1) {$s$};
\node[dnode, below right=9mm and 1mm of x] (t1) {$t$};
\node[dnode, below left=9mm and 1mm of y] (s2) {$s$};
\node[dnode, below right=9mm and 1mm of y] (t2) {$t$};
\draw[dedge] (f) -- node[dlab,left]{$f_{x}$} (x);
\draw[dedge] (f) -- node[dlab,right]{$f_{y}$} (y);
\draw[dedge] (x) -- node[dlab,left]{$x_{s}$} (s1);
\draw[dedge] (x) -- node[dlab,right]{$x_{t}$} (t1);
\draw[dedge] (y) -- node[dlab,left]{$y_{s}$} (s2);
\draw[dedge] (y) -- node[dlab,right]{$y_{t}$} (t2);
\end{tikzpicture}
\end{center}

There are exactly two paths from $s$ to $f$: through $x$, contributing $f_{x}x_{s}$, and through $y$, contributing $f_{y}y_{s}$. Hence $f_{s}=f_{x}x_{s}+f_{y}y_{s}$, and by the same count $f_{t}=f_{x}x_{t}+f_{y}y_{t}$. Reading the tree is the technique. Counting the leaves labelled $s$ tells you how many terms $\pdv{f}{s}$ has, before you compute anything, which is the check that catches the omission.

\tech{The chain rule as a sum over paths}{core}
\cue $z=f(x,y)$ with $x=x(t)$ and $y=y(t)$; or $f\bigl(u(s,t),v(s,t)\bigr)$; or any composite whose inner map has more than one component.
\method Draw the dependency tree from the output down to the independent variables. For each independent variable, sum over all paths reaching it, multiplying the derivative on each edge. With one independent variable the derivatives at the bottom are total and the ones at the top are partial:
\[
\dv{z}{t}=\pdv{f}{x}\dv{x}{t}+\pdv{f}{y}\dv{y}{t},
\qquad
\pdv{z}{s}=f_{u}u_{s}+f_{v}v_{s}.
\]
The one-variable chain rule and the related-rates rule are the cases where the tree has a single branch.

\begin{worked}
Take $z=x^{2}y$ along the unit circle, $x=\cos t$, $y=\sin t$. The tree has two paths, so
\[
\dv{z}{t}=2xy\cdot(-\sin t)+x^{2}\cdot\cos t=-2\cos t\sin^{2}t+\cos^{3}t=\cos t\bigl(\cos^{2}t-2\sin^{2}t\bigr).
\]
At $t=1.1$ this reads $-0.6272108$; differentiating $\cos^{2}t\sin t$ directly and evaluating gives $-0.6272108$.

Now two independent variables. Let $z=\eu^{x}\sin y$ with $x=st$ and $y=s^{2}-t^{2}$. The tree is the one drawn above, so each of $z_{s}$ and $z_{t}$ has two terms:
\[
z_{s}=\eu^{x}\sin y\cdot t+\eu^{x}\cos y\cdot 2s,\qquad
z_{t}=\eu^{x}\sin y\cdot s+\eu^{x}\cos y\cdot(-2t).
\]
At $(s,t)=(1,2)$ we have $x=2$, $y=-3$, and $z_{s}=\eu^{2}\bigl(2\sin(-3)+2\cos(-3)\bigr)=-16.71571$, $z_{t}=\eu^{2}\bigl(\sin(-3)-4\cos(-3)\bigr)=28.21770$. Central differences on $\eu^{st}\sin(s^{2}-t^{2})$ return $-16.71571$ and $28.21770$.

Finally the case that generates the notational trap, where a variable appears both as an independent variable and inside the intermediates. Let $w=f(x,y,t)=xy+t^{2}x$ with $x=\cos t$ and $y=\eu^{t}$. The tree now has three edges from $w$, and $t$ appears as a leaf under $x$, under $y$, and directly under $w$: three paths, three terms.
\[
\dv{w}{t}=f_{x}\dv{x}{t}+f_{y}\dv{y}{t}+f_{t}
=(y+t^{2})(-\sin t)+x\,\eu^{t}+2tx .
\]
At $t=0.7$ this is $0.9980203$, and differencing $w(t)=\eu^{t}\cos t+t^{2}\cos t$ directly gives $0.9980203$. The single term $f_{t}=2tx$ evaluates to $1.0707791$ on its own, a different number entirely. Both are legitimate derivatives of the same function with respect to the same letter; only one of them is the rate of change of $w$ along the path.
\end{worked}

\begin{trap}
Writing one term where the tree has two. This is the path-of-least-resistance error and it is silent: the answer looks like a chain rule, it is dimensionally sensible, and it is wrong. The second trap is notational. If $w=f(x,y,t)$ with $x=x(t)$ and $y=y(t)$, then $\pdv{f}{t}$ means the derivative through the explicit $t$ slot only, while $\dv{w}{t}=f_{x}x'+f_{y}y'+f_{t}$ collects every path. The two are different numbers written with confusingly similar symbols, and a problem that uses both is testing exactly this.
\end{trap}

\begin{seealsob}Partial differentiation with the others frozen; the gradient and directional derivatives; the total differential.\end{seealsob}

\subsection{Second derivatives through the tree}

The tree does not stop at first order. To compute $z_{ss}$ you differentiate $z_{s}=f_{x}x_{s}+f_{y}y_{s}$ again, and the essential observation is that $f_{x}$ and $f_{y}$ are themselves composite functions of $s$ and $t$, so each of them must be run back through the same tree. The result is a product rule on every term, and the count doubles: a first derivative with two paths produces a second derivative with four second-order terms plus two first-order ones. Nearly every mistake at second order is a dropped $x_{ss}$, and it is dropped because the solver differentiated the coefficients and forgot the factors.

The canonical instance is the Laplacian in polar coordinates, worth doing once in full because every other second-order change of variable follows the same pattern. With $x=r\cos\theta$ and $y=r\sin\theta$,
\[
z_{r}=z_{x}\cos\theta+z_{y}\sin\theta ,
\]
and differentiating again in $r$ requires sending $z_{x}$ and $z_{y}$ back down the tree, while $\cos\theta$ and $\sin\theta$ are constants in $r$:
\[
z_{rr}=z_{xx}\cos^{2}\theta+2z_{xy}\sin\theta\cos\theta+z_{yy}\sin^{2}\theta .
\]
In $\theta$ the coefficients are not constant, so the product rule contributes the extra terms:
\[
z_{\theta}=-r\,z_{x}\sin\theta+r\,z_{y}\cos\theta,\qquad
z_{\theta\theta}=r^{2}\bigl(z_{xx}\sin^{2}\theta-2z_{xy}\sin\theta\cos\theta+z_{yy}\cos^{2}\theta\bigr)-r\,z_{r}.
\]
Adding $z_{rr}$ to $r^{-2}z_{\theta\theta}$ makes the cross terms cancel and the squares combine, and the leftover $-r^{-1}z_{r}$ is what forces the middle term of the standard formula:
\[
z_{xx}+z_{yy}=z_{rr}+\frac1r z_{r}+\frac1{r^{2}}z_{\theta\theta}.
\]
For $z=x^{2}y$ at $r=1.7$, $\theta=0.9$, the polar side evaluates numerically to $2.663312$ and the Cartesian side is $2y=2r\sin\theta=2.663311$.

\begin{keynote}[The missing term at second order]
If the inner functions are not linear, $z_{ss}$ contains $f_{x}x_{ss}+f_{y}y_{ss}$ in addition to the four second-order terms. In the polar case $x_{rr}=y_{rr}=0$, which is why the $r$-derivative is clean, and $x_{\theta\theta}=-x$, $y_{\theta\theta}=-y$, which is exactly where the $-r\,z_{r}$ comes from. Whenever a second-order change of variables produces a formula with a first-derivative term in it, that term is the second derivative of the substitution, and its absence from your working is the diagnostic that you differentiated the coefficients only.
\end{keynote}

The tree also explains why implicit differentiation works, why the total differential has the shape it has, and why the gradient is a useful object at all. Each of the next three techniques is the chain rule with a particular tree drawn in advance, which is why they can be reduced to formulas worth memorising. Take direction first.

\section{Direction, and the one object that encodes it}

Once you can differentiate along the two coordinate directions, the natural question is what happens along an arbitrary one. Fix a point $P$ and a unit vector $\hat{\mathbf v}$, and walk from $P$ along $\hat{\mathbf v}$ at unit speed: the position is $P+h\hat{\mathbf v}$, and the rate of change of $f$ is a one-variable derivative in $h$. The tree for that composite has one path per coordinate, so the answer is $f_{x}v_{1}+f_{y}v_{2}+f_{z}v_{3}$, which is a dot product. Collecting the partials into a single vector is therefore not a convenience of notation: it is the statement that all directional information at a point is carried by three numbers.

\tech{The gradient and directional derivatives}{medium}
\cue ``Rate of change in the direction of $\mathbf v$''; ``direction of steepest ascent''; ``normal to the level surface''; ``tangent plane at $P$''.
\method Assemble $\nabla f=(f_{x},f_{y},f_{z})$ at the point. Then
\[
D_{\mathbf v}f=\nabla f\cdot\hat{\mathbf v},\qquad \hat{\mathbf v}=\frac{\mathbf v}{\norm{\mathbf v}} .
\]
Three consequences answer almost every question asked about direction. The maximum of $D_{\mathbf v}f$ over unit $\hat{\mathbf v}$ is $\norm{\nabla f}$, attained when $\hat{\mathbf v}$ points along $\nabla f$; the minimum is $-\norm{\nabla f}$, along $-\nabla f$; and $D_{\mathbf v}f=0$ exactly when $\hat{\mathbf v}\perp\nabla f$, which is why $\nabla f$ is normal to the level set through the point and why the tangent plane to $f=c$ at $P$ is $\nabla f(P)\cdot(\mathbf r-P)=0$.

\begin{worked}
A clean case first. For $f=x^{2}+2y^{2}+3z^{2}$ at $(1,1,1)$, $\nabla f=(2x,4y,6z)=(2,4,6)$. The steepest ascent is along $(1,2,3)$ with rate $\norm{\nabla f}=\sqrt{56}=2\sqrt{14}\approx7.48331$, and along $\mathbf v=(1,2,2)$, whose norm is $3$, the rate is $\tfrac{2+8+12}{3}=\tfrac{22}{3}$. The tangent plane to the ellipsoid $f=6$ at that point is $2(x-1)+4(y-1)+6(z-1)=0$.

A case that exercises the frozen-variable trap. Let
\[
f(x,y,z)=x^{3}z+1+\sin(xy)-x^{zy}
\]
at $(3,0,1)$, along $\mathbf v=(2,4,3)$. The last term is the dangerous one, and it must be differentiated three different ways. In $x$ the power rule applies: $\partial_{x}\bigl(x^{zy}\bigr)=zy\,x^{zy-1}$, which vanishes at $y=0$. In $y$ and in $z$ the exponent is the variable, so each partial carries a factor $\ln x$: $\partial_{y}\bigl(x^{zy}\bigr)=x^{zy}\ln x\cdot z$ and $\partial_{z}\bigl(x^{zy}\bigr)=x^{zy}\ln x\cdot y$. Assembling,
\[
f_{x}=3x^{2}z+y\cos(xy)-zy\,x^{zy-1}=27,\quad
f_{y}=x\cos(xy)-z\,x^{zy}\ln x=3-\ln3,\quad
f_{z}=x^{3}-y\,x^{zy}\ln x=27 .
\]
With $\norm{\mathbf v}=\sqrt{29}$,
\[
D_{\mathbf v}f=\frac{2(27)+4(3-\ln3)+3(27)}{\sqrt{29}}=\frac{147-4\ln3}{\sqrt{29}}\approx 26.48119 ,
\]
and a central difference along the unit vector returns $26.48119$.

The level-surface questions, all answered by the same vector. Take $F(x,y,z)=xyz$ and the surface $F=6$ through $(1,2,3)$. Then $\nabla F=(yz,xz,xy)=(6,3,2)$ at that point, and $\norm{\nabla F}=\sqrt{36+9+4}=7$. From that one vector: the tangent plane is $6(x-1)+3(y-2)+2(z-3)=0$, that is $6x+3y+2z=18$; the normal line is $(1,2,3)+t(6,3,2)$; the steepest increase of $F$ is at rate $7$ along $(6,3,2)$; the directions in which $F$ does not change to first order are exactly those perpendicular to $(6,3,2)$, a two-parameter family spanning the tangent plane; and the rate along any curve lying inside the surface is $0$, since such a curve has $\mathbf r'\perp\nabla F$. Five distinct questions, one computation.
\end{worked}

\begin{trap}
Not normalising $\mathbf v$. This is the commonest slip in the chapter and it produces an answer that is off by the factor $\norm{\mathbf v}$, which is invisible unless you check. If the problem says ``in the direction of the vector $\mathbf v$'' you must divide; if it says ``in the direction of the unit vector'' or gives an angle, you must not. The second slip is reporting $\nabla f$ itself as the answer to ``what is the maximum rate of increase''. The direction is $\nabla f$; the rate is $\norm{\nabla f}$.
\end{trap}

\begin{seealsob}The chain rule as a sum over paths; implicit slope by $-F_{x}/F_{y}$; the exactness test.\end{seealsob}

The gradient being normal to level sets is not a separate fact to memorise. It is the previous technique applied to a curve that stays on a level set: if $\mathbf r(t)$ satisfies $f(\mathbf r(t))=c$, differentiating gives $\nabla f\cdot\mathbf r'=0$, so every tangent vector to the level set is orthogonal to the gradient. That one line is also the whole of implicit differentiation, and it deserves its own technique because the resulting formula is worth carrying in your head.

\section{Curves given implicitly}

\tech[Implicit slope by minus Fx over Fy]{Implicit slope by $-F_{x}/F_{y}$}{core}
\cue A relation $F(x,y)=0$ with the slope wanted, especially when $F$ is symmetric, or large, or has $y$ buried in exponents and logarithms so that collecting $y'$ by hand is error-prone.
\method Regard $y$ as a function of $x$ along the curve and differentiate $F(x,y(x))=0$. The tree has two paths, so $F_{x}+F_{y}\,y'=0$, whence
\[
\dv{y}{x}=-\frac{F_{x}}{F_{y}}\qquad\text{wherever }F_{y}\neq0 .
\]
That is the implicit function theorem in one line, and it is worth memorising because it converts an error-prone algebraic manipulation into two independent partial derivatives and a quotient. The three-variable version is identical: for $F(x,y,z)=0$, $\pdv{z}{x}=-\dfrac{F_{x}}{F_{z}}$ and $\pdv{z}{y}=-\dfrac{F_{y}}{F_{z}}$.

\begin{worked}
The folium of Descartes, $x^{3}+y^{3}-6xy=0$, at $(3,3)$. Take $F=x^{3}+y^{3}-6xy$, which indeed vanishes there since $27+27-54=0$. Then $F_{x}=3x^{2}-6y=9$ and $F_{y}=3y^{2}-6x=9$, so $y'=-1$. Tracing the branch numerically and differencing gives $-1.0000000$.

Now a relation where hand-collecting is genuinely unpleasant: $x^{y}=y^{x}$ at $(2,4)$, a point on the curve since $2^{4}=4^{2}=16$. Set $F=x^{y}-y^{x}$. Each term needs the frozen-variable discipline, because $x$ is a base in one term and an exponent in the other:
\[
F_{x}=y\,x^{y-1}-y^{x}\ln y,\qquad F_{y}=x^{y}\ln x-x\,y^{x-1}.
\]
At $(2,4)$: $F_{x}=4\cdot 8-16\ln 4=32-32\ln 2$ and $F_{y}=16\ln 2-2\cdot 4=16\ln2-8$. Hence
\[
\dv{y}{x}=-\frac{32(1-\ln 2)}{8(2\ln 2-1)}=-\frac{4(1-\ln2)}{2\ln2-1}\approx-3.177399 .
\]
Solving $x^{y}=y^{x}$ numerically for $y$ near $4$ at $x=2\pm10^{-5}$ and differencing gives $-3.177399$. Doing this by logarithmic differentiation instead means writing $y\ln x=x\ln y$, differentiating to $y'\ln x+\tfrac{y}{x}=\ln y+\tfrac{x}{y}y'$, and then collecting; the answer is the same, and there are three more places to lose a term.

The three-variable version, for a surface. On $x^{2}+y^{2}+z^{2}=14$ at $(1,2,3)$, take $F=x^{2}+y^{2}+z^{2}-14$, so $F_{x}=2x$, $F_{y}=2y$, $F_{z}=2z$ and
\[
\pdv{z}{x}=-\frac{2x}{2z}=-\frac{x}{z}=-\frac13,\qquad \pdv{z}{y}=-\frac{y}{z}=-\frac23 ,
\]
which agrees with differentiating the explicit branch $z=\sqrt{14-x^{2}-y^{2}}$ but never required solving for it. That matters as soon as the surface cannot be solved for $z$ in closed form, which is the normal case.
\end{worked}

\begin{trap}
Forgetting the minus sign. The formula is a rearrangement of $F_{x}+F_{y}y'=0$, so the sign is forced; if you cannot remember it, rederive it in that one line rather than guessing. The second failure is applying the formula at a singular point, where $F_{x}=F_{y}=0$. The folium has such a point at the origin, where the curve crosses itself and two tangent lines exist; $-F_{x}/F_{y}$ there is $0/0$ and the implicit function theorem does not apply. A quotient of the form $0/0$ in this formula is a signal that the curve is not locally a graph, not an invitation to use L'Hopital.
\end{trap}

\begin{seealsob}The gradient and directional derivatives; the chain rule as a sum over paths; the exactness test.\end{seealsob}

\begin{keynote}[The cyclic relation, and why it is not $+1$]
For a relation $F(x,y,z)=0$ the three implicit partials multiply to $-1$, not to $1$:
\[
\left(\pdv{x}{y}\right)_{\!z}\left(\pdv{y}{z}\right)_{\!x}\left(\pdv{z}{x}\right)_{\!y}
=\left(-\frac{F_{y}}{F_{x}}\right)\left(-\frac{F_{z}}{F_{y}}\right)\left(-\frac{F_{x}}{F_{z}}\right)=-1 .
\]
Each factor contributes one minus sign and the six partials cancel in pairs. On $x^{2}+y^{2}+z^{2}=14$ at $(1,2,3)$ the three factors are $-2$, $-\tfrac32$, $-\tfrac13$, and their product is $-1$; on the ideal gas law $PV=nRT$ the numerical product of $\pdv{P}{V}$, $\pdv{V}{T}$ and $\pdv{T}{P}$ at $V=0.02$, $T=300$ is $-1.0000008$, the discrepancy being finite-difference error. The subscripts are not decoration: each partial is taken with a different variable held fixed, and the identity is false if you hold the same one throughout.
\end{keynote}

The connection to the previous section is worth stating explicitly, because it makes the formula unforgettable. The curve $F=0$ is a level set of $F$, so $\nabla F=(F_{x},F_{y})$ is normal to it. A normal vector $(A,B)$ corresponds to a tangent line of slope $-A/B$. Substituting $(A,B)=(F_{x},F_{y})$ gives $-F_{x}/F_{y}$ immediately, and the condition $F_{y}\neq0$ becomes the geometric statement that the tangent is not vertical.

\subsection{The same slope three ways}

It is worth taking one relation and differentiating it by every available route, because the comparison is the argument for memorising $-F_{x}/F_{y}$ in miniature. Take $x^{y}=y^{x}$ at the point $(2,4)$.

\emph{By the formula.} Set $F=x^{y}-y^{x}$, compute the two partials independently, divide, negate. Each partial is a single application of the frozen-variable rule to a two-term expression, and no equation is ever rearranged. The answer is $-\dfrac{4(1-\ln2)}{2\ln2-1}$.

\emph{By logarithmic differentiation.} Write $y\ln x=x\ln y$ first, which is a genuine simplification, then differentiate with $y=y(x)$:
\[
y'\ln x+\frac{y}{x}=\ln y+\frac{x}{y}\,y' \ \Longrightarrow\ y'=\frac{\ln y-\tfrac{y}{x}}{\ln x-\tfrac{x}{y}}
=\frac{\ln 4-2}{\ln 2-\tfrac12}=-3.177399 .
\]
Same number, and the intermediate expressions are smaller. This route wins when taking logarithms actually simplifies the relation, which here it does.

\emph{By direct implicit differentiation.} Differentiate $x^{y}=y^{x}$ as it stands. Both sides need the general power rule, $\dvop{x}\bigl(u^{v}\bigr)=u^{v}\bigl(\tfrac{v u'}{u}+v'\ln u\bigr)$, and since $x^{y}=y^{x}=16$ the common factor cancels, leaving $\tfrac{y}{x}+y'\ln x=\ln y+\tfrac{x}{y}y'$, which is the same equation as the previous route reached more cheaply.

Three routes, one answer, and the ranking depends on the relation. Use $-F_{x}/F_{y}$ when $F$ is large or symmetric and no simplification presents itself, because it never asks you to collect $y'$ from both sides. Take logarithms first when the relation is a product, a quotient, or a tower. Never differentiate a relation in place while also trying to simplify it; do one, then the other.

\section{Small changes}

So far every technique has produced an exact derivative. The next one uses derivatives to estimate, which is the form the question takes whenever real measurements are involved. The idea is one dimension of linearisation per variable, added.

\tech{The total differential and error propagation}{medium}
\cue A quantity computed from several measured inputs, each carrying a small uncertainty, with the uncertainty in the output wanted; or any question phrased as ``by approximately how much does $f$ change if $x$ increases by $\ldots$''.
\method The linear approximation in several variables is
\[
\dd f=f_{x}\dd x+f_{y}\dd y,\qquad \Delta f\approx f_{x}\Delta x+f_{y}\Delta y .
\]
For a worst-case error bound take absolute values, $\abs{\Delta f}\leq\abs{f_{x}}\abs{\Delta x}+\abs{f_{y}}\abs{\Delta y}$, which assumes the errors conspire. For independent random errors take them in quadrature, $\sqrt{f_{x}^{2}\Delta x^{2}+f_{y}^{2}\Delta y^{2}}$, which assumes they do not. When $f$ is a product or quotient of powers, take logarithms first: $\dd(\ln f)=\dd f/f$ turns the whole computation into a weighted sum of relative errors, with the weights being the exponents.

\begin{worked}
A cylinder, $V=\pi r^{2}h$. Logarithmic differentiation gives $\ln V=\ln\pi+2\ln r+\ln h$, so
\[
\frac{\dd V}{V}=2\,\frac{\dd r}{r}+\frac{\dd h}{h},
\]
and a $1\%$ error in each input gives a worst case of $3\%$. The exact inflation, $(1.01)^{2}(1.01)-1=0.030301$, confirms that the linear estimate is the leading term. In quadrature the answer would be $\sqrt{2^{2}+1^{2}}=2.236\%$ instead, and the difference between $3\%$ and $2.24\%$ is entirely a modelling decision about whether the errors are correlated.

A case where the answer is not obvious in advance. Two resistors in parallel, $R=\dfrac{R_{1}R_{2}}{R_{1}+R_{2}}$, with $R_{1}=100\pm1$ and $R_{2}=200\pm2$, so $1\%$ each. Here
\[
\pdv{R}{R_{1}}=\frac{R_{2}^{2}}{(R_{1}+R_{2})^{2}}=\frac49,\qquad
\pdv{R}{R_{2}}=\frac{R_{1}^{2}}{(R_{1}+R_{2})^{2}}=\frac19 ,
\]
and $R=66.6\overline{6}$. The worst case is $\tfrac49(1)+\tfrac19(2)=\tfrac23$, which is $1.000\%$ of $R$. The relative error does not grow: a parallel combination of two components each known to $1\%$ is itself known to $1\%$, because the two partials are weights that sum to $1$ when scaled by the relative errors. Scaling both resistors by $1.01$ exactly gives $R\mapsto1.01R$, which is the structural reason. In quadrature the figure drops to $0.745\%$.

A case where an exponent is fractional and therefore reduces the error. The period of a pendulum is $T=2\pi\sqrt{L/g}$, so $\ln T=\ln 2\pi+\tfrac12\ln L-\tfrac12\ln g$ and
\[
\frac{\dd T}{T}=\frac12\frac{\dd L}{L}-\frac12\frac{\dd g}{g}.
\]
A length known to $0.5\%$ and a local $g$ known to $0.2\%$ give a worst case of $0.25\%+0.10\%=0.35\%$ in the period. Notice that the minus sign is irrelevant to the worst case, since the errors are bounded in absolute value and either sign is available, but it is decisive if the two errors have a known common cause. Notice also that measuring $g$ from $T$ and $L$ inverts the weights: $\dd g/g=2\,\dd T/T-\dd L/L$, and now the timing error counts double.
\end{worked}

\begin{trap}
Mixing the two error models. Adding absolute values when the problem says the measurements are independent overstates the error by up to a factor of $\sqrt{n}$; using quadrature when the problem asks for a guaranteed bound understates it. Decide which is wanted from the wording, and say which you used. The second trap is dropping the exponent weight: in $V=\pi r^{2}h$ the radius error counts twice, and in $T=2\pi\sqrt{L/g}$ the length error counts half.
\end{trap}

\begin{seealsob}The chain rule as a sum over paths; the gradient and directional derivatives; logarithmic differentiation.\end{seealsob}

\begin{keynote}[The differential, the tangent plane, and the gradient are one object]
The three formulas
\[
\dd f=f_{x}\dd x+f_{y}\dd y,\qquad
z=f(a,b)+f_{x}(a,b)(x-a)+f_{y}(a,b)(y-b),\qquad
\Delta f\approx\nabla f\cdot\Delta\mathbf r
\]
are the same statement written for three audiences. The first is the calculus of small changes, the second is the geometry of the best affine approximation, and the third is the vector form used for directional questions. If you can move between them freely you will never have to remember which one a problem wants, because you can produce any of them from the list of partials at the point. The single fact underneath all three is differentiability: $f(\mathbf a+\mathbf h)=f(\mathbf a)+\nabla f(\mathbf a)\cdot\mathbf h+o(\norm{\mathbf h})$. Existence of the partials alone is not enough for that, which is the one place where the multivariable theory is genuinely harder than the one-variable theory rather than merely bulkier.
\end{keynote}

\section{Exactness, and the bridge out of this part}

Clairaut's theorem says that if a pair of functions are the two partials of a single $f$, they must satisfy one identity. Read forwards, that is a labour-saving device. Read backwards, it is a test: given a pair $(M,N)$, is there any $f$ at all with $f_{x}=M$ and $f_{y}=N$? The necessary condition is immediate, and on a well-behaved domain it is also sufficient. This is the last technique of Part III and the first technique of the part on differential equations, because an equation written as $M\dd x+N\dd y=0$ is solved in one step the moment the test passes.

\tech[Exactness test My equals Nx]{The exactness test $\pdv{M}{y}=\pdv{N}{x}$}{hard}
\cue A pair $(M,N)$ presented as a candidate gradient, a vector field, or a differential form $M\dd x+N\dd y$; or a first-order equation already arranged as $M\dd x+N\dd y=0$; or a potential $f$ requested from its partials.
\method On a simply connected domain, the form is exact, equivalently the field is a gradient, if and only if
\[
\pdv{M}{y}=\pdv{N}{x} .
\]
When it passes, reconstruct $f$ by integrating $M$ in $x$ and repairing the result:
\[
f=\int M\dd x+g(y),\qquad\text{then set } \pdv{f}{y}=N \text{ to determine } g'(y).
\]
The constant of integration is a function of the other variable, and finding it is the whole second half of the work. Any leftover parameters are fixed by point values or by extra conditions.

\begin{worked}
Suppose $\nabla f=\bigl(3x^{2}+by^{2},\ 2bxy+a\bigr)$ with unknown constants $a,b$, subject to $f(0,0)=0$, $f(1,1)=2$, and $f_{xxx}+\tfrac1x f_{yy}=10$. Find $f(1,2)$.

First the test: $M_{y}=2by$ and $N_{x}=2by$, so the pair is exact for every $a$ and $b$, and a potential exists. Integrate $M$ in $x$:
\[
f=x^{3}+bxy^{2}+g(y).
\]
Differentiate in $y$ and match $N$: $f_{y}=2bxy+g'(y)=2bxy+a$, so $g'(y)=a$ and $g(y)=ay+C$. The condition $f(0,0)=0$ gives $C=0$, so $f=x^{3}+bxy^{2}+ay$.

Now the two remaining conditions. $f_{xxx}=6$ and $f_{yy}=2bx$, so $f_{xxx}+\tfrac1x f_{yy}=6+2b=10$, giving $b=2$. Then $f(1,1)=1+b+a=2$ forces $a=-1$. Hence $f=x^{3}+2xy^{2}-y$ and
\[
f(1,2)=1+8-2=7 .
\]
Differencing $f$ numerically confirms $\nabla f(1.3,0.7)=(6.05,\,2.64)$, matching $(3x^{2}+2y^{2},\,4xy-1)$ at that point.

The same machinery solves a differential equation. The equation $(2xy+y^{2})\dd x+(x^{2}+2xy)\dd y=0$ has $M_{y}=2x+2y=N_{x}$, so it is exact; integrating $M$ gives $f=x^{2}y+xy^{2}+g(y)$, and $f_{y}=x^{2}+2xy+g'(y)=N$ forces $g'=0$. The general solution is the level set $x^{2}y+xy^{2}=C$, obtained without ever separating variables or guessing an integrating factor.
\end{worked}

\begin{trap}
Writing ``$+C$'' instead of ``$+g(y)$'' after the $x$-integration. Holding $y$ fixed means the constant of integration may depend on $y$, and discarding that dependence throws away exactly the information the second step is designed to recover. The second trap is the domain hypothesis. On the punctured plane the field
\[
M=\frac{-y}{x^{2}+y^{2}},\qquad N=\frac{x}{x^{2}+y^{2}}
\]
satisfies $M_{y}=N_{x}=\dfrac{y^{2}-x^{2}}{(x^{2}+y^{2})^{2}}$, verified as $-0.2525040$ for both at $(1.3,0.7)$. Yet it has no potential on the punctured plane, because the only candidate is the polar angle, which cannot be made single-valued around the origin. The test is necessary always and sufficient only where the domain has no holes.
\end{trap}

\begin{seealsob}Partial differentiation with the others frozen; the gradient and directional derivatives; implicit slope by $-F_{x}/F_{y}$.\end{seealsob}

\begin{keynote}[One identity, three readings]
$M_{y}=N_{x}$ is Clairaut, an exactness criterion, and a conservative-field criterion, and these are the same statement seen from three directions. Forwards it lets you choose the cheaper order of mixed differentiation. Backwards it decides whether a pair of functions is a gradient. Physically it decides whether the work done by a field is path independent. When you meet integrating factors in Part V, what you are doing is multiplying a form that fails this test by something that makes it pass.
\end{keynote}

\subsection{When the test fails: integrating factors}

A form that fails $M_{y}=N_{x}$ is not a dead end, because exactness is a property of the way an equation is written, not of the equation itself. Multiplying $M\dd x+N\dd y=0$ by any nonzero $\mu(x,y)$ leaves its solutions untouched but changes both partials, and $\mu$ can often be chosen to make the test pass. The searchable case is when the correcting factor depends on one variable only:
\[
\frac{M_{y}-N_{x}}{N}=h(x)\ \Longrightarrow\ \mu=\eu^{\int h(x)\dd x},
\qquad
\frac{N_{x}-M_{y}}{M}=k(y)\ \Longrightarrow\ \mu=\eu^{\int k(y)\dd y}.
\]
Compute the first quotient; if the $y$ cancels, you have $\mu(x)$. If not, compute the second; if the $x$ cancels, you have $\mu(y)$. If neither cancels, the equation needs a different method, and Part V supplies them.

Take $(3xy+y^{2})\dd x+(x^{2}+xy)\dd y=0$. Here $M_{y}=3x+2y$ and $N_{x}=2x+y$, which differ, so the form is not exact. The first quotient is
\[
\frac{M_{y}-N_{x}}{N}=\frac{x+y}{x^{2}+xy}=\frac{x+y}{x(x+y)}=\frac1x ,
\]
free of $y$, so $\mu=\eu^{\int\dd x/x}=x$. Numerically at $(1.3,0.7)$ the quotient reads $0.769231$ against $1/x=0.769231$. Multiplying through gives $(3x^{2}y+xy^{2})\dd x+(x^{3}+x^{2}y)\dd y=0$, whose partials are both $3x^{2}+2xy$, confirmed as $6.890000$ each at that point. Integrating the new $M$ in $x$ gives $f=x^{3}y+\tfrac12x^{2}y^{2}+g(y)$, and matching $f_{y}=x^{3}+x^{2}y$ forces $g'=0$. The solution curves are
\[
x^{3}y+\tfrac12x^{2}y^{2}=C .
\]
Differencing $f$ at $(1.3,0.7)$ returns $f_{x}=4.186000$ and $f_{y}=3.380000$, matching the corrected $M$ and $N$ exactly. The whole of the theory of first-order exact equations is this test plus this repair, which is why the technique is placed at the end of a chapter on partial derivatives rather than at the start of a chapter on differential equations.

\section{Very high mixed partials}

The last technique in this chapter is the two-variable version of a trick from the chapter on higher-order derivatives: when the order requested is absurd, nobody expects you to differentiate. The requested number is a Taylor coefficient in disguise, and the way to get a Taylor coefficient is to expand, not to differentiate.

\tech{Mixed-partial coefficient extraction}{hard}
\cue An enormous mixed partial at a specific point: $\dfrac{\partial^{134}f}{\partial x^{67}\partial y^{67}}$, $\dfrac{\partial^{100}f}{\partial x^{50}\partial y^{50}}$. The order is far too high for repeated differentiation to be intended.
\method At the origin, the identity is
\[
\frac{\partial^{n+m}f}{\partial x^{n}\partial y^{m}}\Big|_{(0,0)}=n!\,m!\,[x^{n}y^{m}]\,f(x,y),
\]
so expand $f$ as a double series and read off the one coefficient you need. Two structural helps. A diagonal series such as $\sum c_{k}(xy)^{k}$ has only one term contributing to a symmetric mixed partial. And when $f$ contains $\ln(x^{2}+y^{2})$, complexify: $\ln(x^{2}+y^{2})=\ln(x+\iu y)+\ln(x-\iu y)$, and each logarithm has an elementary high derivative in both variables, so the point need not be the origin.

\begin{worked}
The diagonal case. For $f=\dfrac{1}{1-xy}=\sum_{k\geq0}(xy)^{k}$, the only term containing $x^{67}y^{67}$ is $k=67$, whose coefficient is $1$. So
\[
\frac{\partial^{134}f}{\partial x^{67}\partial y^{67}}\Big|_{(0,0)}=67!\cdot67!=(67!)^{2}.
\]
The analogue at $n=1$ can be checked directly: a second difference gives $f_{xy}(0,0)=1.0000=(1!)^{2}$.

The complexified case, evaluated away from the origin:
\[
\frac{\partial^{100}}{\partial x^{50}\partial y^{50}}\Bigl[\arctan\tfrac xy+\arctan\tfrac yx+\ln(x^{2}+y^{2})\Bigr]_{(1,1)} .
\]
The two arctangents sum to $\pi/2$ for $x,y>0$, a constant, so they contribute nothing. For the logarithm write $u=x+\iu y$. Then $\partial_{x}^{n}\ln u=(-1)^{n-1}(n-1)!\,u^{-n}$, and differentiating $u^{-n}$ in $y$ brings down a factor $\iu$ each time, giving
\[
\frac{\partial^{n+m}}{\partial x^{n}\partial y^{m}}\ln(x+\iu y)=(-1)^{n+m-1}\iu^{\,m}(n+m-1)!\,u^{-(n+m)} ,
\]
and the conjugate factor gives the same with $\iu$ replaced by $-\iu$ and $u$ by $\bar u$. At $n=m=50$ the sign factors are $(-1)^{99}\iu^{50}=(-1)(-1)=1$ and likewise for the conjugate, so the value is
\[
99!\bigl[(1+\iu)^{-100}+(1-\iu)^{-100}\bigr].
\]
Since $(1+\iu)^{2}=2\iu$, we get $(1+\iu)^{100}=(2\iu)^{50}=2^{50}\iu^{50}=-2^{50}$, and the same for the conjugate. Hence the answer is
\[
99!\Bigl(-\frac{1}{2^{50}}-\frac{1}{2^{50}}\Bigr)=-\frac{99!}{2^{49}}\approx-1.6578\times10^{141}.
\]
The general formula was checked against direct difference quotients at $n=m=1,2,3$, where it returns $-1$, $-3$ and $-30$ against numerical values $-1.0000$, $-3.0000$ and $-30.0000$.

Two more double series worth knowing, because they cover the non-diagonal cases. First,
\[
\frac{1}{1-x-y}=\sum_{k\geq0}(x+y)^{k}=\sum_{n,m\geq0}\binom{n+m}{n}x^{n}y^{m},
\]
so $\dfrac{\partial^{n+m}f}{\partial x^{n}\partial y^{m}}\big|_{(0,0)}=n!\,m!\,\dbinom{n+m}{n}=(n+m)!$, a striking collapse: the answer depends on $n$ and $m$ only through their sum. Difference quotients at $(n,m)=(1,1),(2,1),(2,2),(3,2)$ return $2.001$, $6.006$, $24.05$ and $120.4$ against $2$, $6$, $24$ and $120$. Second, $\sin(xy)$ is diagonal but with a parity gap: since $\sin u=u-\tfrac{u^{3}}{6}+\tfrac{u^{5}}{120}-\cdots$, the coefficient of $x^{n}y^{n}$ is the coefficient of $u^{n}$ in $\sin u$, which is $0$ for even $n$ and $(-1)^{(n-1)/2}/n!$ for odd $n$. Hence
\[
\frac{\partial^{2n}}{\partial x^{n}\partial y^{n}}\sin(xy)\Big|_{(0,0)}=
\begin{cases}
0, & n \text{ even},\\[2pt]
(-1)^{(n-1)/2}\,n!, & n \text{ odd},
\end{cases}
\]
one factorial rather than two, because the series coefficient already carries a $1/n!$. Numerically at $n=1,2,3$ the difference quotients give $1.000$, $0.000$ and $-6.000$, against $1$, $0$ and $-6$.
\end{worked}

\begin{trap}
Forgetting that there are two factorials. The coefficient formula carries $n!\,m!$, not $(n+m)!$ and not $n!$ alone, and the error is invisible in the symmetric case only until you check magnitudes. The second waste is expanding the double series far past the single monomial you need: decide $(n,m)$ first, then find that one coefficient.
\end{trap}

\begin{seealsob}Maclaurin coefficient extraction; partial differentiation with the others frozen; the annihilated mixed partial.\end{seealsob}

Everything in this chapter has been the one-variable theory with more bookkeeping. The partial derivative is an ordinary derivative with the other letters frozen. The chain rule is the one-variable chain rule summed over the paths of a tree. The gradient is the list of partials, and every directional question is a dot product with it. The implicit formula is the chain rule solved for $y'$. The total differential is the tangent line, one dimension at a time. Exactness is Clairaut read backwards. And the high mixed partial is a Taylor coefficient, exactly as in one variable. If you find yourself doing something in this chapter that has no one-variable ancestor, you have probably made a bookkeeping mistake.

\section{Recognition matrix}
\begin{recog}{@{}p{0.30\textwidth}p{0.36\textwidth}p{0.28\textwidth}@{}}
\rowcolor{mist}\mh{If you see} & \mh{Reach for} & \mh{If that fails} \\
\midrule
\endhead
$f_{x}$ of an expression with $y$ in it & Freeze $y$; rewrite it as a literal constant first & Rename $y$ to $c$ and redo \\
A variable appearing as base and exponent & Three different rules for the three partials & Take logarithms and differentiate \\
$f_{xy}$ requested & Clairaut: pick the order that kills more terms & Check continuity if the function is piecewise \\
$\dfrac{\partial^{n+m}f}{\partial x^{n}\partial y^{m}}$ with $n,m$ large & Coefficient extraction, $n!\,m!\,[x^{n}y^{m}]f$ & Complexify $\ln(x^{2}+y^{2})$ \\
A composite with a multi-component inner map & Draw the dependency tree; sum over paths & Count leaves to check the term count \\
Both $\pdv{f}{t}$ and $\dv{f}{t}$ in one problem & Explicit slot versus all paths; they differ & Write the tree and label the $t$ leaves \\
``In the direction of $\mathbf v$'' & $\nabla f\cdot\mathbf v/\norm{\mathbf v}$ & Check whether $\mathbf v$ was already a unit vector \\
``Steepest ascent'', ``maximum rate'' & Direction $\nabla f$, rate $\norm{\nabla f}$ & Steepest descent is $-\nabla f$ \\
``Normal to the level surface'', tangent plane & $\nabla f(P)\cdot(\mathbf r-P)=0$ & Rewrite the surface as $F=0$ first \\
$F(x,y)=0$, slope wanted & $\dv{y}{x}=-F_{x}/F_{y}$ & If $F_{x}=F_{y}=0$, the point is singular \\
$F(x,y,z)=0$, $\pdv{z}{x}$ wanted & $-F_{x}/F_{z}$ & Solve for $z$ explicitly if possible \\
Measured inputs with tolerances & $\dd f=f_{x}\dd x+f_{y}\dd y$; absolute values for worst case & Quadrature if the errors are independent \\
A product or quotient of powers & Logarithms: relative errors weighted by exponents & Differentiate directly and divide by $f$ \\
A pair $(M,N)$, potential wanted & Test $M_{y}=N_{x}$, then integrate and repair & Check the domain is simply connected \\
$M\dd x+N\dd y=0$ & Exactness test; solution is the level set $f=C$ & Look for an integrating factor \\
$M_{y}=N_{x}$ but no potential exists & The domain has a hole; suspect $\arctan(y/x)$ & Restrict to a simply connected piece \\
\end{recog}
